\chapter{Swimmer s ear}
\index{Swimmer s ear}

\medskip
\noindent\textit{Educational reference; always seek professional advice for individual care.}

\section*{Definition and Overview}
Swimmer's ear is the common name for otitis externa, an infection and inflammation of the outer ear canal, the tube that runs from the opening of the ear to the eardrum. It is called swimmer's ear because it so often follows swimming, when water trapped in the ear canal keeps the skin of the canal wet, soft, and vulnerable to infection. The canal is lined with delicate skin that is normally protected by a thin layer of wax and a mildly acidic environment; when this protection is disrupted, bacteria or fungi that live on the skin multiply and cause infection. The classic features are an itchy, painful ear that hurts when the outer part of the ear is gently pulled, and sometimes discharge or a feeling of fullness. It differs from otitis media, the middle ear infection of childhood, which sits behind the eardrum and needs different treatment. Swimmer's ear is almost always minor and resolves quickly with ear drops, but can become severe if neglected. This chapter explains who gets swimmer's ear, why it happens, how it is recognised, treated, and prevented, and covers the uncommon situations in which it becomes dangerous. Prompt treatment brings rapid relief, and simple habits after swimming prevent most recurrences.

\section*{Epidemiology}
Otitis externa is one of the most common ear problems seen in general practice, and it is more common in warm, humid climates and during the summer swimming season, which is why it features so prominently in local health care. It affects people of all ages, but is particularly common in children and young adults who swim regularly, and in people in hot, dusty environments or who wear hearing aids for long periods. The most important environmental risk factor is moisture: water retained in the canal after swimming, bathing, or showering softens the skin and creates a warm, damp breeding ground for bacteria. The most frequent cause is the bacterium Pseudomonas aeruginosa, which thrives in water, followed by Staphylococcus aureus and other skin bacteria; fungal infections account for a minority of cases and are more likely after prolonged antibiotic ear drops. It is usually a unilateral condition---one ear---but both ears can be affected. Overall, it is a common nuisance rather than a serious illness for most people, but it recurs frequently in individuals who are prone to it, and recurrent or severe disease deserves attention rather than repeated self-treatment.

\section*{Aetiology and Pathophysiology}
The ear canal is a warm, dark, slightly moist tube, and its skin is a barrier that keeps the environment out. Swimmer's ear develops when that barrier fails. The chain usually begins with water remaining in the canal after swimming, which softens and swells the skin and washes away the protective wax and the skin's natural acid layer. Once the skin is wet and alkaline, the bacteria that normally live on it multiply. Any injury, such as scratching with a cotton bud or fingernail, gives the organisms a doorway into the deeper layers, and infection follows.

\begin{itemize}
    \item \textbf{Water exposure:} trapped water softens the skin and removes its protective coating, the essential first step in most cases
    \item \textbf{Trauma:} cotton buds, fingernails, and keys scratch the canal and break the skin barrier; a very common contributing factor
    \item \textbf{Eczema and dermatitis:} skin conditions affecting the canal make the skin fragile and infection-prone
    \item \textbf{Hearing aids, earplugs, and headphones:} these hold moisture and warmth against the canal
    \item \textbf{Excessive cleaning:} over-zealous earwax removal strips the protective layer and irritates the lining
    \item \textbf{Hot, humid climates:} sweating and humidity favour the growth of bacteria and fungi
    \item \textbf{Fungal overgrowth:} after repeated antibiotic drops, fungi can cause a persistent, itching otitis externa
\end{itemize}

The infection causes the lining of the canal to swell, which narrows the tube and can trap pus and debris; it is this swelling against the bony canal wall, which cannot stretch, that makes the pain of swimmer's ear so intense. Left untreated, infection can spread to the ear drum, the cartilage of the outer ear, the nearby lymph nodes, and, in rare and serious cases, to the bone at the base of the skull, a condition called malignant (or necrotising) otitis externa.

\section*{Clinical Features}
The symptoms of swimmer's ear typically build over a few days and are quite distinctive.

\begin{itemize}
    \item \textbf{Itching:} often the earliest symptom, before pain begins
    \item \textbf{Ear pain:} gradually increasing, worse at night, and dramatically worsened by pulling the ear lobe or pressing the tragus, the small bump in front of the ear
    \item \textbf{Feeling of fullness or blockage:} as the canal swells, the ear feels blocked and hearing is slightly muffled
    \item \textbf{Discharge:} a watery, cloudy, or yellowish discharge may appear; fungal infections produce a fluffier material
    \item \textbf{Pain when chewing:} jaw movement moves the ear canal, which hurts when it is inflamed
    \item \textbf{Swelling and redness:} the opening of the ear canal may look red and swollen, and the outer ear may be tender
    \item \textbf{Mild hearing loss:} caused by swelling and debris filling the canal rather than damage to hearing
    \item \textbf{Fever:} usually absent or low-grade; a high fever suggests spread beyond the canal
\end{itemize}

A key point is the distinction from middle ear infection: in swimmer's ear, pulling the ear is very painful, whereas in middle ear infection it is not, because the infection sits behind the eardrum. Children with middle ear infection also more often have fever, earache without canal tenderness, and a recent cold.

\section*{Differential Diagnosis}
Several conditions can mimic swimmer's ear, and the treatment is different for each, which is why a proper look inside the ear matters.

\begin{itemize}
    \item \textbf{Middle ear infection (otitis media):} pain behind the eardrum, usually with a cold, fever, and no pain on pulling the ear; the eardrum itself is red and bulging
    \item \textbf{Furuncle (boil) in the canal:} a localised hair follicle infection producing a pinpoint, very tender swelling rather than generalised canal inflammation
    \item \textbf{Foreign body:} a small object, common in children, can cause discharge and pain from irritation and infection, and is visible on examination
    \item \textbf{Impacted earwax:} a blocked sensation and muffled hearing without the itching, pain, or discharge of infection
    \item \textbf{Eczema or psoriasis of the canal:} itching and scaling without significant pain, discharge, or tenderness
    \item \textbf{Ramsay Hunt syndrome:} severe ear pain with a blistering rash and sometimes facial weakness, caused by the shingles virus
    \item \textbf{Malignant otitis externa:} persistent severe ear pain with discharge in an older person or someone with diabetes, where infection has reached the skull bone
\end{itemize}

\section*{When to Seek Medical Care}
Most cases of swimmer's ear can be treated by a GP, and the sooner treatment begins, the quicker it settles. See a doctor if ear pain is present, if there is discharge, if the ear feels blocked for more than a few days, or if symptoms do not improve after a day or two of basic care. People with diabetes, a weakened immune system, or a previous ear operation such as a cholesteatoma should have ear symptoms assessed early, because their infections can behave more seriously. Do not poke anything into the ear, and do not use over-the-counter drops if the eardrum might be damaged, since some harm the middle ear.

\textbf{Contact your local emergency services for:}
\begin{itemize}
    \item Severe pain with swelling spreading over the cheek, or an ear pushed out by swelling behind it
    \item High fever, confusion, or severe headache
    \item Facial weakness, or difficulty moving the eyes or jaw
    \item Collapse or signs of a severe spreading infection
\end{itemize}

\section*{Diagnosis and Investigations}
Swimmer's ear is usually diagnosed by a doctor looking into the ear with an otoscope, and no further tests are needed for typical cases.

\begin{itemize}
    \item \textbf{Otoscopy:} the doctor examines the canal and eardrum, looking for redness, swelling, discharge, and an intact eardrum, which determines which drops are safe
    \item \textbf{Pain on ear movement:} gentle pulling of the ear or pressing the tragus, which is the hallmark sign of canal infection
    \item \textbf{Swab of discharge:} sent for culture if the infection is severe, recurrent, or failing to respond, to identify the organism and its antibiotic sensitivities
    \item \textbf{Ear clearance:} the doctor may gently suction away debris and pus so drops can reach the skin; this is one of the most effective early steps
    \item \textbf{Blood tests:} rarely needed, but used if a serious spreading infection or malignant otitis externa is suspected, including inflammatory markers and diabetes checks
    \item \textbf{CT or MRI scan:} reserved for suspected malignant otitis externa, to look for infection in the skull bone
    \item \textbf{Hearing assessment:} audiometry may be arranged if hearing loss persists after the infection clears
\end{itemize}

\section*{Management}
The key to treatment is getting antibiotic drops into the ear canal and keeping the canal dry. Most cases clear within about a week.

\begin{itemize}
    \item \textbf{Ear drops:} the mainstay of treatment; combined antibiotic and corticosteroid drops are usually used and are effective against the bacteria that cause swimmer's ear
    \item \textbf{Keeping the ear dry:} avoiding swimming and keeping water out while showering, using a shower cap or a cotton ball smeared with petroleum jelly at the ear opening
    \item \textbf{Suction cleaning:} professionally clearing debris and pus from the canal dramatically improves how well drops work
    \item \textbf{An ear wick:} if the canal is so swollen that drops cannot reach the skin, a small sponge wick is placed to carry the medicine in until the swelling settles
    \item \textbf{Pain relief:} regular simple analgesics such as paracetamol or ibuprofen are often needed for the first few days
    \item \textbf{Oral antibiotics:} needed in a minority of severe cases where infection has spread beyond the canal, or when immunity is weak
    \item \textbf{Fungal treatment:} if the infection is fungal, antifungal drops or creams are used instead of antibiotics
    \item \textbf{Treating the underlying skin problem:} eczema and dermatitis of the canal are treated alongside the infection to prevent recurrence
    \item \textbf{Severe disease:} malignant otitis externa requires hospital treatment with intravenous antibiotics and specialist ear, nose, and throat (ENT) care
\end{itemize}

\section*{Complications}
\begin{itemize}
    \item \textbf{Recurrent otitis externa:} repeated episodes, common in regular swimmers, can leave the canal skin thickened and itchy
    \item \textbf{Stenosis of the canal:} long-standing severe inflammation can narrow the canal and impair hearing
    \item \textbf{Spread of infection:} infection can spread to the cartilage of the outer ear, the ear drum, or the lymph nodes of the neck
    \item \textbf{Malignant (necrotising) otitis externa:} a rare but serious infection of the skull bone, most often in older people with diabetes or weakened immunity, requiring intensive treatment
    \item \textbf{Secondary middle ear infection:} rarely, infection spreads through a damaged eardrum into the middle ear
\end{itemize}

\section*{Prognosis and Outcomes}
With appropriate treatment, the outlook for swimmer's ear is excellent. Most people notice improvement within two to three days of starting drops, and the infection has usually cleared within one to two weeks. Pain is often the slowest symptom to settle, and persistent or worsening pain should prompt a return visit rather than more of the same drops. The main problem for many sufferers is recurrence: people who swim often and remain prone to moisture in the ears may have repeated episodes unless they change their habits after swimming. Malignant otitis externa is a serious infection requiring prolonged intravenous antibiotics and ENT care, but with early treatment the outlook is good. In general, swimmer's ear is preventable and easily treated, and a few simple habits after swimming keep most people free of it.

\section*{Prevention and Self-Care}
\begin{itemize}
    \item After swimming, tilt each ear down to let water drain, or blot the opening with the corner of a dry towel---never poke anything inside
    \item Keep the ears dry with a shower cap while showering if you are prone to infections
    \item Do not clean the ear canals with cotton buds, fingers, or any object; wax is protective and the canal cleans itself
    \item A few drops of a commercially available acidic "swimmer's drops" solution in each ear after swimming can help prevent infection
    \item If you use hearing aids or earplugs, keep them clean and give your ears regular breaks
    \item Treat underlying eczema or dermatitis of the ear with your usual skin treatment, and see a GP if it is not controlled
    \item See a doctor early for ear pain or discharge rather than waiting for it to worsen
\end{itemize}

\section*{Sources and Further Reading}
The following organisations provide reliable information about ear infections, swimming, and ear care.

\begin{itemize}
    \item healthdirect Australia. \textit{Swimmer's ear (otitis externa).} https://www.healthdirect.gov.au/
    \item NHS. \textit{Outer ear infection (otitis externa).} https://www.nhs.uk/
    \item Mayo Clinic. \textit{Swimmer's ear: symptoms, causes, and treatment.} https://www.mayoclinic.org/
    \item Australian Department of Health and Aged Care. \textit{Ear health and hearing resources.} https://www.health.gov.au/
    \item Royal Children's Hospital Melbourne. \textit{Ear infections in children: a parent's guide.} https://www.rch.org.au/
    \item World Health Organization (WHO). \textit{Ear and hearing care information.} https://www.who.int/
\end{itemize}

\clearpage
